\section{Future work}
\label{sec:future}



Right now, our Level-2 selection currently focuses on minimizing the score while maximizing \texttt{output\_seq\_len}. In future work, we can extend it by also considering \texttt{input\_seq\_len}. Optimizing only the score and \texttt{output\_seq\_len} can favor configurations that get good results simply by using a longer input window, to have more context. However, another model can reach a similar score and horizon with a shorter \texttt{input\_seq\_len}, which is preferable because it reduces memory and performance. So, we can think like this for the model selection: (i) minimize score, (ii) maximize \texttt{output\_seq\_len}, and (iii) minimize \texttt{input\_seq\_len}, for example, using a three-dimensional Pareto front.


Autoregressive decoding and scheduled sampling are both step-by-step: each prediction depends on the output generated at the previous step. This dependency limits parallelization and makes both training and inference slow, especially for long horizons.

Scheduled sampling has also some limitations:
\begin{itemize}
  \item It \textbf{does not fully remove exposure bias}. This means that errors can still get worse when the model must rely on its own predictions.

  \item It also \textbf{creates a mismatch between train and test}, because training mixes ground truth inputs and model generated inputs in a way that is different from inference.
  
  \item In some cases it can be unstable and return suboptimal or unpredictable behaviour.
  
\end{itemize}

To reduce the step-wise limitation, a practical direction is to use more parallelizable variants. 
For Transformers, a simple option is \emph{two-pass scheduled sampling}, which applies scheduled sampling through a two pass decoding strategy while keeping most computation parallel \cite{mihaylova2019scheduledtransformers}.